% !TEX root = ../Projektdokumentation.tex
\section{Qualitätssicherung}
\label{sec:Qualitaetssicherung}
Die Qualitätssicherung erfolgte durch automatisierte und manuelle Tests. Mit
Pytest wurden vor allem Validierungen, Benutzerverwaltung, Audit-Funktionen und
Formularlogik geprüft. Zusätzlich wurden die wichtigsten Arbeitsabläufe direkt
über die Weboberfläche getestet, darunter Anmeldung, Rollenvergabe,
Mandantentrennung sowie die Pflege von Mail- und IP-Adressen.

Diese Mischung war für das Projekt passend, weil sie sowohl die fachlichen
Regeln als auch die Bedienung abdeckt. Die automatisierten Tests helfen dabei,
Fehler bei späteren Änderungen schneller zu erkennen. Die manuellen Tests
waren zusätzlich notwendig, um Sichtbarkeiten, Navigation und Fehlermeldungen
im Zusammenspiel der einzelnen Bereiche zu prüfen. Ein Auszug aus dem
automatisierten Zugriffstest befindet sich im \Anhang{app:Testauszug}.

\begin{table}[htbp]
\centering
\begin{tabular}{|p{4.0cm}|p{4.0cm}|p{4.0cm}|c|}
\hline
\textbf{Testfall} & \textbf{Erwartet} & \textbf{Tatsächlich} & \textbf{Check} \\
\hline
Admin legt Benutzer an & Benutzer wird gespeichert und zur Benutzerliste weitergeleitet & HTTP 303 auf \Code{/users/page}, Benutzer in Datenbank vorhanden & OK \\
\hline
Rolle und Firmenzuordnung & Neuer Benutzer besitzt Rolle \Code{User} und die richtige Firmen-ID & Rolle und Firmen-ID entsprechen den Testdaten & OK \\
\hline
Passwortspeicherung & Passwort wird nicht im Klartext gespeichert, sondern als Hash geprüft & \Code{verify\_password} bestätigt den gespeicherten Hash & OK \\
\hline
User öffnet Benutzerverwaltung & Zugriff wird verweigert & HTTP 403 & OK \\
\hline
User legt weiteren Benutzer an & Zugriff wird verweigert & HTTP 403 & OK \\
\hline
\end{tabular}
\caption{Auszug der automatisierten Zugriffstests}
\label{tab:Testauszug}
\end{table}

\subsection{Soll-Ist-Vergleich von Leistung, Zeit und Inhalt}
\label{sec:SollIstVergleich}
Die geplanten Funktionen wurden erreicht. Die Anwendung ist
mehrmandantenfähig, stellt ein rollenbasiertes Berechtigungskonzept bereit und
ermöglicht die zentrale Pflege von Firmen, Benutzern, Mail-Adressen
und IP-Adressen. Zusätzlich wurden Audit-Log, Login-Schutz und Rollback für
ausgewählte Änderungen umgesetzt.

\begin{table}[htbp]
\centering
\begin{tabular}{|p{5.8cm}|p{2.3cm}|p{2.8cm}|}
\hline
\textbf{Aspekt} & \textbf{Soll} & \textbf{Ist} \\
\hline
Mandantenfähige Verwaltung & umgesetzt & umgesetzt \\
\hline
Rollen- und Rechtesystem & umgesetzt & umgesetzt \\
\hline
Auditierbare Änderungen & umgesetzt & umgesetzt \\
\hline
Rollback einzelner Änderungen & optional & umgesetzt \\
\hline
Projektzeit & 80 h & eingehalten \\
\hline
\end{tabular}
\caption{Soll-Ist-Vergleich}
\label{tab:SollIst}
\end{table}

Inhaltlich gab es keine grundlegenden Abweichungen vom Projektziel. Im
Projektverlauf wurde etwas mehr Wert auf Sicherheit und Nachvollziehbarkeit
gelegt. Das war sinnvoll, weil bei einer zentralen Verwaltungsanwendung klar
sein muss, wer welche Daten bearbeiten darf und welche Änderungen später
nachvollzogen werden können.

\subsection{Übergabe und Abnahme}
\label{sec:UebergabeUndAbnahme}
Nach Abschluss der Implementierung wurde die Anwendung dem Auftraggeber
vorgeführt. Dabei wurden die zentralen Abläufe gemeinsam durchgegangen:
Anmeldung, Pflege von Firmen, Benutzern, Mail-Adressen und IP-Adressen,
Rollenprüfung sowie Nachvollziehbarkeit über das Audit-Log. Die Rückmeldungen
aus dem Review waren positiv, da die bisher manuelle Pflege, nachvollziehbarer
und einheitlicher abgebildet wurde.

Das Projektergebnis wurde fachlich abgenommen. Offene Punkte betrafen keine
fehlenden Kernfunktionen, sondern mögliche spätere Erweiterungen für einen
produktiven Ausbau.

\subsection{Bewertung der Testergebnisse}
\label{sec:BewertungTestergebnisse}
Die durchgeführten Tests bestätigten, dass die wesentlichen Pflegeabläufe
stabil funktionieren. Besonders wichtig war dabei die Kombination aus Login,
Rollenprüfung, Mandantentrennung und Auditierung. Offene Punkte betreffen keine
fehlenden Hauptfunktionen, sondern mögliche spätere Erweiterungen wie
weitergehende Integrations- oder Lasttests für einen produktionsnahen Ausbau.
